\section*{Chapter \ref{ch:mx:tick-size-queues-and-priority} --- Tick Size, Queues and Priority}

\begin{solution}{exo:mx:tick-size-queues-and-priority:1}
$0.01/4=25$ basis points and $0.01/400=0.25$ basis points. The four-dollar stock is more likely to be large-tick, but the answer depends on the spread it
would have without the constraint, which depends on its volatility and activity: a very active, low-volatility stock can be large-tick at a high price.
\end{solution}

\begin{solution}{exo:mx:tick-size-queues-and-priority:2}
$\hat\eta=40/(2\times160)=0.125$; implicit spread $2\times0.125\times1=0.25$ cents. Yes: $\eta$ is well below one half.
\end{solution}

\begin{solution}{exo:mx:tick-size-queues-and-priority:3}
$P_2=\tfrac{0.2}{0.25}\cdot\tfrac{0.25}{0.30}\cdot\tfrac{0.30}{0.35}=0.571$. Expected time: $1/0.2+1/0.25+1/0.3=12.33$ seconds.
\end{solution}

\begin{solution}{exo:mx:tick-size-queues-and-priority:4}
Front: $0.7\times(2.5-0.4)=1.47$ cents per share; back: $0.06\times(2.5-0.9)=0.096$. The front is worth 1.37 cents per share more: USD 6.87 for 500
shares.
\end{solution}

\begin{solution}{exo:mx:tick-size-queues-and-priority:5}
Before--after: $9\,960-2\,400=7\,560$. \omterm{def:mx:tick-size-queues-and-priority:did}{Difference-in-differences}: $7\,560-(2\,350-2\,750)=7\,960$. It is causal if, without the new grid, the treated
stocks' depth would have changed like the controls' (parallel trends); here it holds by construction, because the simulator draws the stocks alike.
\end{solution}

\begin{solution}{exo:mx:tick-size-queues-and-priority:6}
The three windows of one stock share its liquidity and its random draws, so they are positively correlated and carry less information than three
independent observations; the independent formula ignores that and understates the error by a factor of 1.7. Report the clustered one; with only 16
clusters, a small-sample correction or a wild bootstrap is safer still.
\end{solution}

\begin{solution}{exo:mx:tick-size-queues-and-priority:7}
Two hours per rule on seeds 7 and 8: under time priority 0.70 of front orders and 0.05 of those behind 20 or more receive a fill, worth 1.28 and 0.081 cents
per posted share; under pro rata 0.78 and 0.58, worth 0.72 and 0.093. Pro rata gives each resting order a slice of every fill in proportion to its size,
so an order deep in a long queue still receives only a small part of its size, and the value per posted share falls with the length of the queue it
shares.
\end{solution}

\begin{solution}{exo:mx:tick-size-queues-and-priority:8}
Depth at the best rises mechanically when several one-cent levels are pooled at one five-cent price: the same shares are counted at fewer prices. The
SEC staff found no significant change in depth within five cents of the best for the first two groups, and \omterm{def:mx:decomposing-the-spread:quoted}{quoted spreads} rose five to seven basis
points against the control group. The simulated experiment shows the same: depth at the best grows 3.8-fold on a five-cent grid with no new orders at
all, while every \omterm{def:mx:the-limit-order-book:orders}{market order} pays 4.5 times more.
\end{solution}

\begin{solution}{pb:mx:tick-size-queues-and-priority:1}
\textbf{1.} One whose spread is one tick almost always; the tick must be compared with the spread the asset would have unconstrained, which depends on
volatility and activity.
\textbf{2.} It holds \texttt{firm.tape}'s order flow fixed and changes only the price grid (limit prices rounded away from the market).
\textbf{3.} 1.07 cents and 94.2\%; 5.03 cents and 99.4\%.
\textbf{4.} 1\,740 to 6\,560 shares, 8.7 to 33.6 orders.
\textbf{5.} 0.56 and 2.53 cents of effective half-spread: the same \omterm{def:mx:the-limit-order-book:orders}{market orders} cross a spread that is now one five-cent tick.
\textbf{6.} Expected profit per share from joining; $V(n)=P_n(c-a(n))$.
\textbf{7.} 11.3, 24.1, 55.7 and 81.4 seconds.
\textbf{8.} 0.76, 0.44 and 0.13 predicted at 0, 3 and 20 orders ahead against 0.54, 0.22 and 0.06: the model ignores new levels in front and faster
cancellations when the price moves.
\textbf{9.} Front 0.57, 0.66, 1.26 and 2.05 cents per posted share; back 0.001, 0.005, 0.094 and 0.188.
\textbf{10.} The back fills only when the flow has eaten the whole queue, which is more often flow that keeps moving the price: at one cent, $-0.71$ cents
at the front against 0.48 behind 20 orders.
\textbf{11.} $\hat\eta=N_c/(2N_a)$ over one-tick changes of the traded price.
\textbf{12.} 0.108, 0.058, 0.019 and 0.007; 0.22, 0.23, 0.19 and 0.14 cents.
\textbf{13.} The implicit spread is close to a property of the asset: the \omterm{def:mx:decomposing-the-spread:quoted}{quoted spread} follows the grid, the asset's own cost of a price change does not.
\textbf{14.} NMS stocks priced at USD 1.00 or more whose time-weighted average \omterm{def:mx:decomposing-the-spread:quoted}{quoted spread} is USD 0.015 or less, for quotes and orders.
\textbf{15.} About 400 stocks each: quoting in five cents; quoting and trading in five cents; the same plus trade-at.
\textbf{16.} Spreads rose for every group between the two periods, control included.
\textbf{17.} \omterm{def:mx:decomposing-the-spread:quoted}{Quoted spreads} up five to seven basis points against the control; depth at the best up, mostly mechanically; depth within five cents not
significantly changed in the first two groups.
\textbf{18.} $+7\,960$ shares; 2\,270 clustered by stock, 1\,340 unclustered.
\textbf{19.} \emph{Named result}: multiplying the tick by five on the same flow multiplies the spread by 4.7 (1.07 to 5.03 cents), the displayed size at the
best by 3.8, the median time to fill by 4.9 and the value of the front of the queue by 2.2 (0.57 to 1.26 cents per posted share); the pilot measured a
five-to-seven basis point rise of \omterm{def:mx:decomposing-the-spread:quoted}{quoted spreads} against controls and a mostly mechanical rise of depth at the best.
\textbf{20.} Liquidity providers at the front of the queue, paid by every order that crosses the wider spread.
\end{solution}

\begin{solution}{iq:mx:tick-size-queues-and-priority:1}
Its spread is one tick almost always and its queues are long; trading it is a question of queue position and timing (joining early, not paying the
spread), while a small-tick stock has a sparse book, spread changes and competition on price.
\iqlookfor{Spread pinned at one tick; queues; price against time competition.}
\end{solution}

\begin{solution}{iq:mx:tick-size-queues-and-priority:2}
The front fills more often and against less informed flow; the back fills rarely and mostly when the price is moving through it. It is worth more when the
tick is large relative to the asset's natural spread, queues are long, and priority is by time.
\iqlookfor{Both factors (probability and conditional adverse selection); the tick; the priority rule.}
\end{solution}

\begin{solution}{iq:mx:tick-size-queues-and-priority:3}
A birth--death queue: \omterm{def:mx:the-limit-order-book:orders}{market orders} at rate $\mu$ take the front, orders ahead cancel at $\theta$, ours at $\nu$: $P_n=\tfrac{\mu}{\mu+\nu}\prod_k\tfrac{\mu+k
\theta}{\mu+k\theta+\nu}$. It leaves out price moves (a better level forming, the level emptying), state-dependent cancellations, market-order sizes and
\omterm{def:mx:order-types-and-their-uses:hidden}{hidden orders}.
\iqlookfor{The formula; its missing price dynamics.}
\end{solution}

\begin{solution}{iq:mx:tick-size-queues-and-priority:4}
\omterm{def:mx:tick-size-queues-and-priority:did}{Difference-in-differences}: regress spreads on treated, after and their interaction with stock and date controls, standard errors clustered by stock;
check pre-trends and the balance of the groups, and report by pre-period spread (the constraint binds only where the old spread was below the new tick).
\iqlookfor{Control group; interaction; clustering; heterogeneity.}
\end{solution}

\begin{solution}{iq:mx:tick-size-queues-and-priority:5}
Position matters less and size more: fills are split by size, so I post more than I want to trade and manage the overfill risk, watch the top-order and
minimum-allocation rules, and cancel faster when the queue's total size could fill me too much at once.
\iqlookfor{Oversizing; allocation rules; fill risk.}
\end{solution}

\begin{solution}{iq:mx:tick-size-queues-and-priority:6}
Check the share of time the spread is one tick and the queue lengths, then estimate $\eta$ from one-tick continuations and alternations: $\eta$ well below
one half, with a spread pinned at one tick, means the tick is large relative to the asset's implicit spread $2\eta\delta_{\mathrm{tick}}$.
\iqlookfor{One-tick share; queue lengths; the uncertainty-zone estimator.}
\end{solution}
